\documentclass[11pt,a4paper,nofootinbib]{article}
\usepackage[utf8]{inputenc}
\usepackage[T1]{fontenc}
\usepackage[german]{babel}
\usepackage{amsmath}
\usepackage{amsfonts}
\usepackage{amssymb}
\usepackage{bm}
\usepackage{geometry}
\usepackage{booktabs}
\usepackage{cite}
\usepackage{hyperref}

\geometry{a4paper, margin=25mm}

\hypersetup{
	colorlinks=true,
	linkcolor=blue,
	citecolor=blue,
	filecolor=magenta,      
	urlcolor=cyan,
	pdftitle={Effektive arithmetisch-geometrische Dualitaet},
	pdfpagemode=FullScreen,
}

\title{\textbf{Effektive arithmetisch-geometrische Dualität im quaternionischen Zahlenraum}\\
	\large --- Eine effektive Feldtheorie auf dem Hurwitz-Gitter (\#Energiedoku) ---}
\author{\textbf{Thomas Hoffbauer} \\
	\small Bamberg, Deutschland}
\date{\small 29. Mai 2026}

\begin{document}
	
	\maketitle
	
	\begin{abstract}
		In dieser Arbeit formulieren wir ein rigoroses mathematisches Framework, das die diskrete algebraische Faktorisierung natürlicher Zahlen über eine dynamische Dreikörper-Traktorie auf dem Hurwitz-Quaternionen-Gitter in eine kontinuierliche, pseudo-riemannsche Feldtheorie überführt. Wir zeigen, dass die aus dem mikroskopischen Gitterprozess gewonnene Uhrenpräzision $\eta$ als direkte Quelle für den makroskopischen Energie-Impuls-Tensor dient. Dies erlaubt die Herleitung einer arithmetischen Schwarzschild-Metrik, deren Krümmungsskala direkt durch die quantenhafte informationelle Stabilität des zugrundeliegenden Zahlenraums gesteuert wird. Numerische Analysen des deviatorischen Momententensors und der Zustandsübergänge (SLERP) verifizieren die strikte geometrische Rigidität des Innenraums.
	\end{abstract}
	
	\section{Einführung und kanonische Minkowski-Einbettung}
	Ein zentrales Anliegen des Bamberg-Modells (\#Energiedoku) ist die geometrische Rigideisierung des Zahlenraums zur Untersuchung arithmetischer Interferenzen. Jede natürliche Zahl $n \in \mathbb{N}$ lässt sich kanonisch anhand ihrer Primfaktorzerlegung in drei strukturell disjunkte Klassen unterteilen:
	\begin{enumerate}
		\item \textbf{Das Fundament (2,3-Plattenanteil):} $n_{2,3} = 2^{\alpha} \cdot 3^{\beta}$
		\item \textbf{Die Halbgruppe (5,7-Halb-Kleinscher Anteil):} $n_{5,7} = 5^{\gamma} \cdot 7^{\delta}$
		\item \textbf{Der Kern (Ganz-Kleinscher Anteil):} $n_{>7} = \prod_{p_i > 7} p_i^{\epsilon_i}$
	\end{enumerate}
	
	Wir definieren eine Abbildung $\Psi: \mathbb{N} \to \mathbb{M}^{1,3}$ in den Minkowski-Raum mit der metrischen Signatur $(+,-,-,-)$. Die raumzeitlichen Koordinaten $x^\mu = (x^0, x^1, x^2, x^3)$ werden durch die logarithmischen Amplituden dieser arithmetischen Klassen besetzt, um die multiplikative Struktur in ein additives Vektorfeld zu überführen:
	\begin{align}
		x^0(n) &= c \cdot \ln(n_{2,3}) \quad \text{(Chronologische Zeitachse)} \\
		x^1(n) &= \ln(n_{5,7}) \\
		x^2(n) &= \ln(n_{>7}) \\
		x^3(n) &= \theta(n)
	\end{align}
	Hierbei repräsentiert $x^3(n) = \theta(n)$ einen phasen- und winkelabhängigen Term, welcher die relative Lage der Primzahlzyklen beschreibt.
	
	\section{Das geschlossene Dreikörperproblem der Zahlkomponenten}
	Innerhalb der Einbettung fassen wir die drei räumlichen Koordinatenkomponenten als wechselwirkende Punktmassen $m_1, m_2, m_3$ auf, deren Trägheit proportional zu ihrer arithmetischen Komplexität skaliert. Der gravitative Schwerpunkt $R^\mu$ dieses abgeschlossenen Systems entwickelt sich entlang der Eigenzeit $\tau$ (welche starr mit der $x^0$-Achse gekoppelt ist) nach:
	\begin{equation}
		R^\mu(\tau) = \frac{\sum_{i=1}^3 m_i x_i^\mu(\tau)}{\sum_{i=1}^3 m_i}
	\end{equation}
	Die Evolution dieser Komponenten umeinander beschreibt eine deterministische Traktorie $\gamma_n(\tau)$ im Minkowski-Raum. Die Überlagerung dieser Einzeltraktorien bei der Interferenz mehrerer Zahlen generiert ein kontinuierliches Hintergrundfeld mit dem kollektiven Energie-Impuls-Tensor:
	\begin{equation}
		T_{\mu\nu} = \sum_{i} T_{\mu\nu}^{(n_i)}
	\end{equation}
	
	\section{Die Einsteinschen Feldgleichungen der Arithmetik}
	Die kollektive Dichte der Primzahlzyklen krümmt die asymptotisch flache Raumzeit. Die Geometrie (beschrieben durch den metrischen Tensor $g_{\mu\nu}$) folgt den Einsteinschen Feldgleichungen:
	\begin{equation}
		G_{\mu\nu} + \Lambda g_{\mu\nu} = \kappa_{\text{eff}} T_{\mu\nu}^{(\text{Arith})}
	\end{equation}
	Für den Außenbereich isolierter Resonanzen setzen wir ein statisch-sphärisches Modell an. Die Geodätenbewegung einer Testzahl $m$ folgt der klassischen Gleichung:
	\begin{equation}
		\frac{d^2 x^\mu}{d\tau^2} + \Gamma^\mu_{\alpha\beta} \frac{dx^\alpha}{d\tau} \frac{dx^\beta}{d\tau} = 0
	\end{equation}
	Die Christoffel-Symbole $\Gamma^\mu_{\alpha\beta}$ kodieren hierbei die Scheinkräfte der arithmetischen Verteilung. Numerische Simulationen in \texttt{SageMath 10.5} liefern für den Außenraum ($r > r_A$) die exakten nicht-verschwindenden Komponenten:
	\begin{align}
		\Gamma^t_{t,r} &= \frac{r_A}{2(r^2 - r \cdot r_A)} \\
		\Gamma^r_{t,t} &= \frac{r \cdot r_A - r_A^2}{2r^3} \\
		\Gamma^r_{r,r} &= -\frac{r_A}{2(r^2 - r \cdot r_A)}
	\end{align}
	
	\section{Hauptsatz: Effektive arithmetisch-geometrische Dualität}
	Das zentrale Ergebnis dieses Formalismus lässt sich in folgendem Theorem zusammenfassen:
	
	\medskip
	\noindent\textbf{Satz (Effektive arithmetisch-geometrische Dualität).}\\
	\textit{Sei $\eta = \frac{T}{S} > 0$ die aus einem abgeschlossenen Hurwitz-Gitterprozess gewonnene Uhrenpräzision (wobei $T$ die Anzahl der Quanten-Ticks und $S$ die generierte Gitter-Entropie bezeichnet). Definiere den arithmetischen Horizontradius $r_A = \alpha \eta$ mit dem Skalenfaktor $\alpha > 0$ und die effektive arithmetische Energiedichte $\rho_{\text{Arith}} = \hbar_{\text{eff}} \eta$.\\
		Dann existiert ein effektives, statisch-sphärisches Modell $(\mathcal{M}, g)$ mit
		\begin{equation}
			G_{\mu\nu} + \Lambda g_{\mu\nu} = \kappa_{\text{eff}} T_{\mu\nu}^{(\text{Arith})},
		\end{equation}
		wobei der Energie-Impuls-Tensor stückweise definiert ist als:
		\begin{equation}
			T_{\mu\nu}^{(\text{Arith})} = \begin{cases} 
				\rho_{\text{Arith}} u_\mu u_\nu & \text{für } r \le r_A \quad \text{\small(Quellträger / Innenraum)} \\ 
				0 & \text{für } r > r_A \quad \text{\small(Außenraum)} 
			\end{cases}
		\end{equation}
		Im Außenbereich ist die Geodätenbewegung durch die zugehörigen Christoffel-Symbole eindeutig bestimmt; die mikroskopische Größe $\eta$ steuert die makroskopische Krümmungsskala $r_A$.}\\
	\medskip
	
	\section{Numerische Verifikation und Gitter-Scherung}
	Zur empirischen Validierung des Modells wurden iterative Gitter-Transformationen unter einer Störungsamplitude von $0.18$ ausgewertet. Die numerische Quantenuhr liefert stabile zeitliche Kennzahlen:
	\begin{equation}
		T = 16, \quad S = 1.8096 \implies \eta \approx 8.8415 \text{ Ticks/Entropie}
	\end{equation}
	
	Die kontinuierliche Festkörper- und Scherungsanalyse des Innenraums ($r \le r_A$) zeigt eine akkumulierte Scherenergie (Frobenius-Norm) von $0.001962$. Die isotrope Volumendehnung (seismische Spur) beträgt $\text{Sp}(T_{\mu\nu}) = -3.003200$, was einer strikt kompressiven raumzeitlichen Bindung entspricht. Der zugehörige deviatorische Momententensor für den reinen Scherungsanteil bestimmt sich zu:
	\begin{equation}
		M_{\text{dev}} = \begin{pmatrix} 
			0.04026667 & 0.4992 & 0.52 \\ 
			0.4992 & -0.03973333 & 0.48 \\ 
			0.52 & 0.48 & -0.00053333 
		\end{pmatrix}
	\end{equation}
	Da die Spur dieses Deviators verschwindet ($\sum M_{ii} \approx 0$), bleibt das System frei von unphysikalischen Verkippungen. 
	
	Der direkte quantitative Vergleich der Zustandsübergänge bei der Mitte des Phasenpfads ($t = 0.5$) beweist die Notwendigkeit der geodätischen Pfadwahrung: Während eine lineare Interpolation (LERP) eine Norm von $0.8660$ erzeugt und somit über einen Volumenkollaps mechanischen Stress induziert, garantiert das sphärisch-lineare Thermostat (EABC, Grenzwert $0.01$) eine \textbf{SLERP-Norm von exakt 1.0000}. Die perfekte geometrische Rigidität auf der Einheits-Quaternionenkugel bleibt invariant gewahrt.
	
	\section{Fazit und Ausblick}
	Das formulierte Modell zeigt, dass informationelle Masse als rein topologisches Phänomen der Gitterstabilität interpretiert werden kann. Zukünftige Arbeiten im Rahmen des \textit{Bamberger Protokolls} werden untersuchen, wie die Nullstellen der Riemannschen Zeta-Funktion als stabile Orbits oder Photonensphären auf der kritischen Resonanzlinie $r = 3 r_A$ lokalisiert werden können.
	
\end{document}
